% ===========================================================================
% SECTION 1 — INTRODUCTION
% Target: 800-1000 words
% ===========================================================================

\section{Introduction}
\label{sec:introduction}

Wind energy is one of the fastest-growing renewable energy sources, and
the continued growth of installed capacity places increasing demands on
turbine control systems \citep{pao2009}. Above the rated wind speed
(Region~III), a variable-speed, variable-pitch wind turbine must regulate
rotor speed to its rated value by actuating blade pitch, while the
generator torque is held at (or switched into) its rated value
\citep{bianchi2006}. This regulation problem is nonlinear, subject to
actuator rate and range constraints, and driven by a highly stochastic
disturbance (turbulent wind); Model Predictive Control (MPC) is
well-suited to it precisely because it handles input and state
constraints explicitly and can anticipate disturbance effects over a
prediction horizon \citep{rawlings2017}, in contrast to classical
gain-scheduled PI pitch
controllers \citep{bianchi2006,pao2009,laks2009}. Realizing these
specific advantages in closed loop requires exercising them --
multiple simultaneously active constraints, or genuine disturbance
preview -- neither of which the single-constraint, no-preview
protocol used throughout this paper does; a gain-scheduled PI
reference is included in Section~\ref{sec:closed_loop_results}
accordingly, and this paper's central comparison is instead among
surrogate architectures' fidelity to the true-model MPC they replace,
not a claim that MPC outperforms classical control under this
specific protocol. The practical obstacle to
nonlinear MPC (NMPC) for this problem is computational: at every control
step, the solver must repeatedly evaluate the internal plant model and,
for gradient-based solvers, its Jacobian which is expensive when that
model is the full nonlinear aerodynamic and structural dynamics of the
turbine \citep{schlipf2013,evans2022}.

\subsection{ML-MPC: State of the Art}
\label{subsec:soa}

A natural response to this computational bottleneck is to replace the
physical plant model inside the MPC prediction step with a learned
surrogate that is cheaper to evaluate. \citet{klein2024} demonstrated a Local Linear Neuro-Fuzzy Model
(LLNFM) surrogate for wind turbine load estimation  embedded within an
Extended Kalman Filter state estimator (not an MPC internal predictor),
validated in simulation only  the direct reference and starting point
for the LLNFM implementation adopted in this work
(Section~\ref{sec:system_modeling}). A related paper from the same
research group, \citet{klein2025}, separately reports the first
full-scale experimental validation of MPC on an operating 3\,MW turbine
(3~hours of continuous field operation)  but using a robust
linear-time-varying MPC combined with an Extended Kalman filter, with no
learned surrogate of any kind. We cite both papers as motivation and
precedent for, respectively, the LLNFM architecture and the field-testing
standard this line of work aspires to, without implying either paper
combines both properties simultaneously. Beyond neuro-fuzzy
surrogates, Gaussian Process (GP) regression has an established track
record as an MPC-internal model precisely because it returns a calibrated
predictive uncertainty alongside its point prediction, enabling cautious
or uncertainty-aware control formulations \citep{kocijan2004,hewing2020} —
a property we exploit directly in the GP-based cost of
Section~\ref{subsec:gp_uncertainty_cost}. Within wind energy specifically,
\citet{liu_wind_2018} use a Mat\'ern-kernel GP to predict short-term wind
speed for preview-based load mitigation, but  unlike the present
work  as an external disturbance forecast feeding the controller
rather than as the MPC's internal dynamics model. Data-driven
reinforcement learning has separately been applied to wind turbine
torque/pitch control \citep{xie_data-driven_2023}, offering a
model-free alternative to the surrogate-based MPC approach studied
here; we return to this comparison in
Section~\ref{sec:discussion}. Physics-Informed Neural Networks
(PINNs), which embed a governing differential equation as a soft
constraint on the training loss, have become a general framework for
combining data-driven flexibility with physical consistency
\citep{raissi2019,karniadakis2021}, and Temporal Convolutional Networks
(TCNs) have been shown to be a competitive, often superior alternative to
recurrent architectures for generic sequence modeling tasks
\citep{bai2018}. Each of these four paradigms neuro-fuzzy, GP, PINN,
and convolutional/recurrent sequence models has therefore been
independently validated as a plausible MPC surrogate architecture in some
application domain. Beyond ML surrogates specifically, the broader
engineering practice of MPC continues to evolve rapidly
\citep{schwenzer2021}, and linear-quadratic-optimal alternatives to
nonlinear MPC for wind turbine power tracking, without a learned
internal model, have also recently been proposed and validated in
OpenFAST \citep{grapentin2022,grapentin2024}  a useful point of
comparison for the computational-cost trade-offs discussed in
Section~\ref{sec:mpc_framework}, though outside this paper's
ML-surrogate scope.

What is missing, to our knowledge, is a systematic comparison of these
paradigms \emph{within a single closed-loop MPC framework, on the same
plant, under the same solver}, paired with an explicit operational
feasibility analysis of running each surrogate inside a real-time SQP
solve rather than only an open-loop accuracy comparison. Prior work
typically evaluates one surrogate family in isolation
\citep{klein2024,kocijan2004,hewing2020}; the interaction between
surrogate architecture and the specific numerical behavior of a
gradient-based NMPC solver which, as we show in
Section~\ref{subsec:operational_feasibility}, can render an
otherwise-accurate surrogate practically unusable is, to our knowledge,
not addressed in the existing wind turbine ML-MPC literature.

\subsection{Contributions}
\label{subsec:contributions}

This paper makes four contributions:

\begin{enumerate}
\item A systematic comparison of six ML surrogate architectures (LLNFM,
LSTM, GP, PINN, TCN, SW-MLP) as internal state predictors for nonlinear
MPC of wind turbine rotor speed, evaluated both open-loop (one-step
prediction accuracy) and closed-loop (tracking performance under
turbulent wind), across two model generations (V1, V2).
\item A five-stage operational feasibility diagnostic that discovered,
contrary to an initial hypothesis of architectural incompatibility,
that deep-learning surrogates can be successfully integrated into
closed-loop SQP-based \texttt{nlmpc} when the repeated \texttt{predict()}
calls inside the solver loop are replaced with a verified manual
forward-pass computation in double precision. The diagnostic traced a
$10^{2}$--$10^{3}\times$ per-call overhead to \texttt{predict()}'s
internal dlarray dispatch and JIT recompilation, not to architecture,
model size, or MATLAB version, and demonstrated that the remedy is
both effective (restoring closed-loop feasibility for all six
surrogates tested) and general (applicable to any \texttt{dlnetwork} or
similar architecture). A second, independent failure mode is
additionally identified and quantified: \texttt{predict()}'s
single-precision output can silently zero out \texttt{fmincon}'s
default finite-difference gradient estimate for a given decision
variable, corrupting solver feasibility independent of any real-time
deadline, with direct evidence from both a controlled precision test
and matched closed-loop infeasibility rates.
\item A Region~II/III generator torque switching correction for the
two-state plant model, calibrated at the turbine's actual rated operating
point rather than its theoretical aerodynamic optimum, without which
long-horizon closed-loop simulation enters an irrecoverable operating
trap a modeling detail we show is not a minor implementation nicety but
a precondition for valid training data and closed-loop evaluation.
\item Applying the same diagnostic standard as Finding 2 to our own
GP-based uncertainty-penalized MPC cost formulation, tracing a
previously reported 'zero pitch activity' result to a frozen-actuator
solver-infeasibility artifact  withdrawn accordingly, and reported
itself as this paper's fourth contribution in place of the original,
now-retracted claim
(Section~\ref{subsec:gp_frozen_actuator_finding}).
\end{enumerate}


\begin{table}[H]
	\centering
	\caption{Prior surrogate-in-controller studies most directly comparable
		to this work.}
	\label{tab:prior_studies}
	\small
	\setlength{\tabcolsep}{3pt}
	\begin{tabularx}{\textwidth}{
			>{\raggedright\arraybackslash}p{1.8cm}
			>{\raggedright\arraybackslash}p{2.1cm}
			>{\raggedright\arraybackslash}X
			>{\raggedright\arraybackslash}p{2.0cm}
			>{\centering\arraybackslash}p{1cm}
			>{\centering\arraybackslash}p{1.5cm}
			>{\raggedright\arraybackslash}X
		}
		\toprule
		Study & Surrogate family & Plant & Solver &
		Closed-loop? & Cost reported? & Comparison scope \\
		\midrule
		
		\citet{witter_creating_2025}
		& LSTM, GRU, CNN, Transformer
		& MBS wind turbine
		& Turbine's own controller (no MPC)
		& Yes
		& No
		& Single turbine, one plant surrogate role \\
		
		\citet{salzmann_real-time_2022}
		& MLP (various sizes)
		& Quadrotor
		& SQP (acados)
		& Yes
		& Yes
		& Robotics, not wind \\
		
		\citet{klein2024}
		& LLNFM
		& NREL 5-MW (sim.)
		& EKF (not MPC)
		& N/A
		& No
		& Load estimation only \\
		
		\citet{klein2025}
		& None (robust LTV)
		& 3\,MW field turbine
		& LTV MPC
		& Yes
		& No
		& Field validation, no surrogate \\
		
		\citet{fu_gaussian_2025}
		& GP
		& NREL 5-MW (floating)
		& MPC
		& Yes
		& Not confirmed
		& Individual pitch, uncertainty-aware cost \\
		
		\citet{liu_wind_2018}
		& GP (wind forecast)
		& Generic turbine
		& MPC
		& Yes
		& No
		& GP for disturbance forecast, not dynamics \\
		
		\textbf{This paper}
		& LLNFM, LSTM, GP, PINN, TCN, SW-MLP (6)
		& NREL 5-MW
		& SQP (\texttt{nlmpc})
		& Yes
		& Yes
		& Systematic multi-architecture, single plant/solver \\
		
		\bottomrule
	\end{tabularx}
	
	\vspace{4pt}
	{\footnotesize
		\textit{No prior study identified compares multiple ML surrogate
			families as an MPC internal predictor, on the same plant and solver,
			with an explicit operational-feasibility diagnostic  the gap this
			paper addresses.}
	}
\end{table}

%% [TODO -- item 4.4, decide with supervisor] Referee requests a minimum
%% of 4 reference categories. Status: 3/4 already covered above
%% (surrogate-in-controller outside wind: salzmann_real-time_2022;
%% GP-in-wind incl. 2025 floating turbine study: fu_gaussian_2025,
%% liu_wind_2018; reference controller for this turbine: jonkman2009,
%% already used for turbine params and the Section~6 PI controller).
%% MISSING: explicit-MPC/neural-approximation-of-control-law
%% literature (category 1) -- no citation yet, needs adding.
%% ALSO UNRESOLVED: fu_gaussian_2025 is cited only once (this table)
%% and never discussed against this paper's own Contribution 4
%% (GP-uncertainty-cost frozen-actuator finding) -- Fu et al.
%% successfully deploy uncertainty-aware GP-MPC while we find ours
%% structurally infeasible; this contrast needs either a deeper read
%% of Fu et al.'s cost formulation to identify the actual difference,
%% or an explicit acknowledgment that the discrepancy is unexplained.
%% Both points pending supervisor decision.

\subsection{Relationship to a Companion Manuscript}
\label{subsec:companion_manuscript}

This paper is an original-research contribution reporting a single
closed-loop comparative study and diagnostic investigation. It shares
partial reference overlap and topical scope with a companion, separately
submitted systematic literature review covering machine-learning
surrogates for model predictive control across both individual wind
turbine and wind farm applications; that review synthesizes a broad
multi-source literature landscape and does not report new closed-loop
simulation results of its own. The two manuscripts are complementary but
distinct in contribution: the companion review situates this paper's
findings within the wider published literature, while this paper
reports original simulation-based results (Sections~\ref{sec:ml_surrogates}\ref{sec:closed_loop_results})
and the \texttt{predict()}-overhead diagnostic
(Section~\ref{subsec:operational_feasibility}) not present in the review.

\subsection{Paper Organization}
\label{subsec:organization}

Section~\ref{sec:system_modeling} presents the plant, torque-switching,
and wind models. Section~\ref{sec:dataset_generation} describes the
training data generation protocol. Section~\ref{sec:ml_surrogates}
details all six surrogate architectures and their open-loop accuracy.
Section~\ref{sec:mpc_framework} presents the NMPC formulation, surrogate
integration, and operational feasibility analysis. Section~\ref{sec:closed_loop_results}
reports closed-loop V1 and V2 results, their comparison, and a Monte Carlo
robustness analysis. Section~\ref{sec:discussion} discusses the main
findings, limitations, and directions for future work, and
Section~\ref{sec:conclusion} concludes.
